\documentclass[10pt,a4paper]{article}
\usepackage{amsmath,amssymb,amsthm,graphicx}
% Card-local minimal preamble (journal unit ahl-197-on-the-existence-of-logarithmic-and-orbifold-jet-2024, P2-ahl16 precedent).
%% Verbatim from the pinned source: the author theorem environments and macro block,
%% so every macro and environment of the counted spans is available.
% Omitted and why: the documentclass line (the journal class sigma is not on the host;
%% the assembler supplies the class) and page-geometry/font blocks (they fight the
%% forced xelatex engine).
\usepackage{amsmath,amssymb,amsthm,mathtools}

%% \usepackage{physics} is deliberately NOT loaded (L60449 incident 2026-10-15):
%% the span's physics-trigger macros \abs and \Tr are both defined by the author
%% block itself (\DeclarePairedDelimiter{\abs} at pinned source L24, \Tr below),
%% and physics's built-ins collide with them (\Tr, \rank) while its starred
%% delimiter machinery (\opbraces) chokes on the span's multi-line \abs* argument
%% that the author's mathtools version accepts.  The author's own paired
%% delimiter definitions are restored verbatim below (keep-pattern regex gap,
%% fixed 2026-10-15: DeclarePairedDelimiter lines are now kept).
%\DeclareMathOperator{\Im}{Im}
\DeclareMathOperator{\lcm}{lcm}
\DeclareMathOperator{\Arg}{Arg}
\DeclareMathOperator{\Hom}{Hom}
\DeclareMathOperator{\SM}{SM}
\DeclareMathOperator{\Herm}{Herm}
\DeclareMathOperator{\Tr}{Tr}
\DeclareMathOperator{\Pic}{Pic}
\DeclareMathOperator{\Proj}{Proj}
\DeclareMathOperator{\Vol}{Vol}
\DeclareMathOperator{\Span}{Span}
\DeclareMathOperator{\rank}{rank}
\DeclareMathOperator{\diff}{d}
\DeclareMathOperator{\dbar}{\bar\partial}
\DeclareMathOperator{\ddbar}{\partial\bar\partial}
%% Author's own paired delimiters, verbatim from pinned source lines 24-25
%% (the keep-pattern regex lacked DeclarePairedDelimiter; fixed 2026-10-15).
\DeclarePairedDelimiter{\abs}{\lvert}{\rvert}
\DeclarePairedDelimiter{\Abs}{\lVert}{\rVert}
\newenvironment{equation'}{
\begin{equation}
\tag{\theequation\ensuremath{'}}}{
\end{equation}
}
\newenvironment{equation_}[1]{
\begin{equation}
\tag{\ensuremath{\theequation_{#1}}}}{
\end{equation}
}
\newenvironment{eqlist}{
\begingroup \leftskip\parindent \rightskip\parindent \noindent
\begin{minipage}{\linewidth-2\parindent}} {
\end{minipage}
\endgroup}


%\let\Im\relax




%\let\Im\relax
\newcommand*\rS{\mathrm{S}}
\newcommand*\Id{\mathrm{Id}}
\newcommand*\Sing{\mathrm{Sing}}
\newcommand*\orb{\mathrm{orb}}
\newcommand*\FS{\mathrm{FS}}
\newcommand*\loc{\mathrm{loc}}
\newcommand{\rma}{\mathrm{a}}
\newcommand{\rmb}{\mathrm{b}}
\newcommand*\cI{\mathcal{I}}
\newcommand*\cJ{\mathcal{J}}
\newcommand*\cO{\mathcal{O}}
\newcommand*\bC{\mathbb{C}}
\newcommand*\bD{\mathbb{D}}
\newcommand*\bE{\mathbb{E}}
\newcommand*\bG{\mathbb{G}}
\newcommand*\bN{\mathbb{N}}
\newcommand*\bP{\mathbb{P}}
\newcommand*\bQ{\mathbb{Q}}
\newcommand*\bR{\mathbb{R}}
\newcommand*\bS{\mathbb{S}}
\newcommand*\bZ{\mathbb{Z}}
\newcommand*\bB{\mathbb{B}}
\newcommand*\bOne{{\mathchoice {\rm 1\mskip-4mu l} {\rm 1\mskip-4mu l} {\rm 1\mskip-4.5mu l}{\rm 1\mskip-5mu l}}}
\newcommand\bDel{\mbox{$\Delta$\newline \scalebox{.6}{\raisebox{3pt}{$\mkern-22mu\not{}\mkern12mu$}}}}
%\newcommand\bDel{\ooalign{$\Delta$\\ \scalebox{.6}{\raisebox{3pt}{\,$\not{}$}}}}
\newcommand\bDelta{\mathchoice{\bDel}{\bDel}{\scalebox{.7}{\bDel}}{\scalebox{.5}{\bDel}}}
%\newcommand\bDelta{\mbox{$\Delta\mkern-2.5mu\not{}$}}
\let\leq\leqslant
\let\geq\geqslant
\newcommand*\bydef{\coloneqq}
\newcommand*{\mk}{\mkern -1mu}
\newcommand*{\Mk}{\mkern -2mu}
\newcommand*{\mK}{\mkern 1mu}
\newcommand*{\MK}{\mkern 2mu}
\newcommand*{\romanenumi}{\renewcommand*{\theenumi}{\roman{enumi}}}
\newcommand*{\Romanenumi}{\renewcommand*{\theenumi}{\Roman{enumi}}}
\newcommand*{\alphenumi}{\renewcommand*{\theenumi}{\alph{enumi}}}
\newcommand*{\Alphenumi}{\renewcommand*{\theenumi}{\Alph{enumi}}}
\let\oldtilde\tilde
\renewcommand*{\tilde}[1]{\mathchoice{\widetilde{#1}}{\widetilde{#1}}{\oldtilde{#1}}{\oldtilde{#1}}}
%\let\oldhat\hat
%\renewcommand*{\hat}[1]{\mathchoice{\widehat{#1}}{\widehat{#1}}{\oldhat{#1}}{\oldhat{#1}}}
\let\oldexists\exists
\renewcommand*{\exists}{\mathrel{\oldexists}}
\let\oldforall\forall
\renewcommand*{\forall}{\mathrel{\oldforall}}
\newtheorem{cdrthm}{Theorem}[section]
\newtheorem{theo}[cdrthm]{Theorem}
\newtheorem{prop}[cdrthm]{Proposition}
\newtheorem{lemm}[cdrthm]{Lemma}
\newtheorem{coro}[cdrthm]{Corollary}
\newtheorem{defi}[cdrthm]{Definition}
\newtheorem{rema}[cdrthm]{Remark}
\newtheorem{conj}[cdrthm]{Conjecture}
\numberwithin{equation}{section}

%% compile hygiene: sub-50pt overfull reports are cosmetic;
\hfuzz=50pt
\title{Lower bound for twisted k-jet sections through an approximate Zariski decomposition}
\author{Campana, Frederic; Darondeau, Lionel; Demailly, Jean-Pierre; Rousseau, Erwan}
\date{doi:10.5802/ahl.197; CC BY 4.0}
\begin{document}
\maketitle
\setcounter{section}{5}
\setcounter{cdrthm}{2}
\noindent\textit{Editorial note.} All delivered mathematical text is the author's own, copied from the pinned journal source file sources/ahl-197-on-the-existence-of-logarithmic-and-orbifold-jet-2024/pinned-original-source.tex (On the existence of logarithmic and orbifold jet differentials); no statement or proof step was rewritten, reordered or removed.\par\medskip
\noindent\textit{Editorial note.} The counted proof is the author's complete argument: Published printed number is Corollary 5.3. The bracketed mention [Fuj94] in the proof is a literal source text token, not a \textbackslash{}cite command: the span carries zero \textbackslash{}cite and its only \textbackslash{}ref is the corollary's own label (self-citation at the end of the proof, resolves internally). All displays inside the span are unnumbered; no labelled equations.\par\medskip
\noindent\textit{Editorial note.} The label coro5.3 names the counted statement, and the proof environment immediately following it in the source is delivered verbatim as its proof.\par\medskip
\noindent\textit{Editorial note.} The counted statement is carried by the author's own coro environment in the article source.\par\medskip
\noindent\textit{Editorial note.} Journal-unit assembly: card-local preamble (author theorem environments and macros verbatim), article class, no bibliography in the counted span (the span carries no \textbackslash{}cite commands); references leaving the span are declared as navigation stubs carrying the article's printed numbers.\par\medskip
\noindent\textit{Editorial scope.} the counted statement and its own proof environment only\par\medskip
\begin{coro}\label{coro5.3}
If $F$ is an arbitrary $\bQ$-line bundle over $X$, one has
\begin{multline*}
h^0\bigg(X_k(V),\cO_{X_k(V)}(m)\otimes\pi_k^*\cO
\Big(-\tfrac{m}{ kr}\Big(1+\tfrac{1}{ 2}+\cdots+\tfrac{1}{ k}\Big)F\Big)\bigg)
\\
\geq \tfrac{m^{n+kr-1}}{(n+kr-1)!}
\tfrac{(\log k)^n}{ n!\,k!^r}\left(
\Vol(K_V\otimes F^{-1})-O\left((\log k)^{-1}\right)\right)-o\left(m^{n+kr-1}\right),
\end{multline*}
$k$-jet differentials of degree $m$ twisted by the appropriate power of $F$ if $K_V\otimes F^{-1}$ is big.
\end{coro}

\begin{proof}
The volume is computed here as usual, i.e.\ after performing a suitable modifi\-cation $\mu:\smash{\widetilde X}\to X$ which converts $K_V$ into an invertible sheaf. There is of course nothing to prove if
$K_V\otimes F^{-1}$ is not big, so we can assume $\Vol(K_V\otimes F^{-1})>0$. Let us fix smooth hermitian metrics $h_0$ on $T_X$ and
$h_F$ on $F$. They induce a metric $\mu^*(\det h_0^{-1}\otimes h_F^{-1})$ on $\mu^*(K_V\otimes F^{-1})$ which, by our definition of $K_V$, is a smooth metric. By the result of Fujita [Fuj94] on approximate Zariski decomposition, for every $\delta>0$, one can find a modification $\mu_\delta:\smash{\widetilde X_\delta}\to X$ dominating
$\mu$ such that
\[
\mu_\delta^*\left(K_V\otimes F^{-1}\right) =\cO_{\widetilde X_\delta}(A+E)
\]
where $A$ and $E$ are $\bQ$-divisors, $A$ ample and $E$ effective, with
\[
\Vol(A)=A^n\geq \Vol\left(K_V\otimes F^{-1}\right)-\delta.
\]
If we take a smooth metric $h_A$ with positive definite curvature form
$\Theta_{A,h_A}$, then we get a singular hermitian metric $h_Ah_E$ on
$\mu_\delta^*(K_V\otimes F)$ with poles along $E$, i.e.\ the quotient
$h_Ah_E/\mu^*(\det h_0^{-1}\otimes h_F)$ is of the form $e^{-\varphi}$ where
$\varphi$ is quasi-psh with log poles $\log\abs{\sigma_E}^2$ (mod $C^\infty(\smash{\widetilde X_\delta}))$ precisely given by the divisor $E$. We then only need to take the singular metric $h$ on $T_X$ defined by
\[
h=h_0e^{\tfrac{1}{ r}(\mu_\delta)^*\varphi}
\]
(the choice of the factor $\tfrac{1}{ r}$ is there to correct adequately the metric on $\det V$). By construction $h$ induces an admissible metric on $V$ and the resulting curvature current $\eta_F=\Theta_{K_V,\det h^*}-\Theta_{F,h_F}$ is such that
\[
\mu_\delta^*\eta_F = \Theta_{A,h_A} +[E],\qquad
\text{$[E]={}$current of integration on $E$.}
\]
Then the $0$-index Morse integral in the complement of the poles is given by
\[
\int_{X(\eta,0)\setminus S}\eta_F^n=\int_{\widetilde X_\delta}\Theta_{A,h_A}^n=A^n\ge
\Vol\left(K_V\otimes F^{-1}\right)-\delta
\]
and Corollary~\ref{coro5.3} follows from the fact that $\delta$ can be taken arbitrary small.
\end{proof}

\end{document}
